\chapter{Logic programming as proof search}\label{ch:horn}

\Cref{ch:intro} gave the overall plan of this document: read Miller and Nadathur's \emph{Programming with
Higher-Order Logic}~\cite{MillerNadathur2012} with a miniKanren programmer's eyes, and find out what
\micro{} would need in order to host its ideas. This chapter lays the first-order foundation. It covers the
book's Introduction and its Chapters~1 and~2, plus the parts of the Teyjus appendix needed to run the
examples. Nothing here goes beyond Prolog in expressive power. What is new, compared with the usual
Prolog textbook, is the \emph{point of view}. A logic program is a set of formulas, a query is a formula,
and running the query means \emph{searching for a proof} in a sequent calculus. Each logical connective
has a fixed search meaning, independent of the program. Atomic goals are solved by \emph{backchaining}
on program clauses. A language counts as a logic programming language when this goal-directed search
is \emph{complete}.

For a miniKanren programmer the punchline comes early. Pure miniKanren, with \scm{=°},
\scm{fresh°}, \scm{cond°} and recursive relations, \emph{is} first-order Horn-clause logic programming, except for its
search strategy and its lack of types. This chapter makes that statement precise, so that later
chapters can say exactly what Horn clauses \emph{lack}: local constants and local clauses
(\cref{ch:hh}), and $\lambda$-terms as data (\cref{ch:lambda,ch:lambdatree}).

\section{Computation as deduction}\label{sec:horn-deduction}

The Introduction distinguishes two ways of using logic in computer science \book{1--3}. In
\emph{computation-as-model}, a computation is a mathematical structure (states, transitions), and
logic talks \emph{about} it from outside, as Hoare triples or temporal modalities do. In
\emph{computation-as-deduction}, the logical objects themselves (formulas, terms, proofs) \emph{are} the
elements of the computation. This second view splits in two again:
\begin{itemize}
  \item \textbf{proof normalization}: the state is a proof (a $\lambda$-term), and computing means reducing
    it to normal form. This is the basis of functional programming;
  \item \textbf{proof search}: the state is a \emph{sequent}, that is, a goal together with its assumptions,
    and computing means searching for a derivation of it. The way sequents change during the search
    \emph{is} the dynamics of the computation. This is the basis of logic programming.
\end{itemize}

Proof search in logic programming is not the proof search of mathematicians. Mathematicians use
\emph{lemmas}, and the proof rule for lemmas is the \emph{cut rule}: to prove $B$, state a lemma $C$, then
prove $C$ and prove $B$ from $C$. Gentzen's cut-elimination theorem~\cite{Gentzen1935} says that cut can
always be eliminated, for classical and for intuitionistic logic. Cut-free proofs are huge and of
little use to a mathematician, but they are exactly what a machine can look for. So in logic
programming ``the cut rule is excluded from the execution of logic programs, and computation is based on
the search for cut-free proofs''; cut and cut elimination remain useful to \emph{reason about} programs
\book{3}. We come back to this in \cref{sec:horn-logic}.

\paragraph{Programs, goals, sequents.} A formula is either part of a \emph{logic program}, which gives
meaning to the non-logical constants it mentions, or a \emph{goal}, which we try to derive from the
program \book{3}. The state of the search is a collection of sequents
\[
  \seq{\Sig}{\Prog}{G}
\]
where $\Sig$ is a \emph{signature} (the typed non-logical constants available), $\Prog$ is the program and
$G$ is the goal. The Introduction states the principle that drives the whole book \book{3--4}:
\begin{itemize}
  \item if $G$ is \emph{not atomic}, the search commits to decomposing its top-level connective. The
    \emph{search semantics} of a connective is fixed and does not depend on the program;
  \item if $G$ is \emph{atomic}, the program is consulted. The search \emph{backchains}: it looks in
    $\Prog$ for an implication whose conclusion matches the atom, and then tries to prove the premise.
\end{itemize}
``Logic programming can be seen abstractly as a logical framework in which a strategy that alternates
between goal reduction and backchaining is complete'' \book{4}.

\paragraph{How can a sequent change?} The Introduction then asks a qualitative question \book{4}: if
proving $\seq{\Sig}{\Prog}{A}$ leads to proving $\seq{\Sig'}{\Prog'}{A'}$, what can differ between the two?
For \emph{Horn clauses}, $\Sig'=\Sig$ and $\Prog'=\Prog$: signature and program are global and fixed, and
only the atom changes. For \emph{hereditary Harrop formulas}, $\Sig'\supseteq\Sig$ and
$\Prog'\supseteq\Prog$: the search can add constants and clauses that are visible only in part of the
computation, which gives modular programming and data abstraction. Since more of the dynamics then
comes from the logic and not from the particular predicates, ``the metatheory of the logic can be of
more value in proving properties of those programs'' \book{4}. This chapter treats the first case; the
second is \cref{ch:hh}.

\section{Typed first-order terms}\label{sec:horn-terms}

Before writing relations we need data. Book Chapter~1 fixes the language of \emph{terms}. It has no
predicates and no programs yet. It uses actual \lp{λProlog} syntax and is the $\lambda$-free special
case of what \cref{ch:lambda} generalises.

\subsection{Sorts, type constructors and kinds}

λProlog has built-in \emph{sorts} (atomic types) \book{10--11}: \lp{int}, \lp{real}, \lp{string},
\lp{in_stream}, \lp{out_stream}, and \lp{o}, the type of formulas, which has a special status
(\cref{sec:horn-formulas}). It also has the built-in \emph{type constructor} \lp{list}, so \lp{(list int)}
is a type. Type constructors are classified by \emph{kinds}:
\[
  \langle\text{kind}\rangle ::= \typ{type} \mid \typ{type}\to\langle\text{kind}\rangle .
\]
A kind is therefore just an arity: \lp{int : type}, \lp{list : type -> type}. New type constructors are
declared with \lp{kind} \book{12}:
\begin{lprolog}
kind  pair  type -> type -> type.
\end{lprolog}
after which \lp{(pair string int)} is a type. Type expressions are type variables (uppercase tokens),
arrows $\tau_1\to\tau_2$ (written \lp{->}, right associative), and type constructors applied to exactly as
many arguments as their arity \book{12--13}.

\begin{definition}[Target type, order {\book{13}}]\label{def:horn-order}
Every type can be written uniquely as $\tau=\tau_1\to\cdots\to\tau_n\to\tau_0$ with $\tau_0$ not an arrow.
$\tau_0$ is the \emph{target type} and $\tau_1,\dots,\tau_n$ are the \emph{argument types}; $\tau$ is
\emph{functional} when $n\ge1$. The \emph{order} of a type is
\[
  \operatorname{ord}(\tau)=0 \text{ if $\tau$ is not functional},\qquad
  \operatorname{ord}(\tau_1\to\tau_2)=\max(\operatorname{ord}(\tau_1)+1,\ \operatorname{ord}(\tau_2)).
\]
\end{definition}
The order counts how deeply the arrow is nested \emph{to the left}. For example \lp{int} and
\lp{list (int -> int)} have order~0 (the order does not look inside constructor arguments),
\lp{int -> string -> o} has order~1, and \lp{int -> (int -> o) -> o} has order~2 \book{14}. In the
first-order language every declared constant has a type of order at most~1 that does not mention
\lp{o}, every variable has a type of order~0, and type variables may only be instantiated with order-0
types \book{13--16}. The last condition is needed because order is not stable under type substitution:
\lp{A -> list A -> o} has order~1, but instantiating \lp{A} with \lp{int -> int} gives a type of
order~2 \book{13--14}.

\subsection{Terms, polymorphism and typing}

All symbols are \emph{constants}; a constant of functional type is a \emph{value constructor}. Lists use
the polymorphic constants \lp{nil : list A} and \lp{:: : A -> list A -> list A}, where \lp{::} is infix and
right associative, so \lp{2 :: 1 :: nil} has type \lp{list int} \book{15}. Constants are declared with
\lp{type} \book{16}:
\begin{lprolog}
type  pr  A -> B -> pair A B.
\end{lprolog}
A crucial asymmetry governs polymorphism \book{17}:
\begin{itemize}
  \item each \emph{occurrence of a constant} may take its own instance of the declared type. In
    \lp{(2 :: 1 :: nil) :: (1 :: nil) :: nil} the first two \lp{nil}s have type \lp{list int} and the last
    one has type \lp{list (list int)};
  \item a \emph{variable} has one type at all its occurrences. So \lp{((X :: nil) :: X :: nil)} cannot be
    typed, because $X$ would need both a type $\tau$ and the type \lp{list} $\tau$.
\end{itemize}
The rules are those of the type assignment calculus of the book's Figure~1.2 \book{18}, with judgments
$\Sig;\Gamma\Vdash_f t:\tau$, where $\Sig$ types the constants and $\Gamma$ the variables:
\begin{mathpar}
  \inferrule{c:\sigma\in\Sig \\ \tau\vartriangleleft_f\sigma}{\Sig;\Gamma\Vdash_f c:\tau}
  \and
  \inferrule{x:\tau\in\Gamma}{\Sig;\Gamma\Vdash_f x:\tau}
  \and
  \inferrule{\Sig;\Gamma\Vdash_f g:\tau_1\to\tau_2 \\ \Sig;\Gamma\Vdash_f t:\tau_1}{\Sig;\Gamma\Vdash_f (g\ t):\tau_2}
\end{mathpar}
Here $\tau\vartriangleleft_f\sigma$ (``$\tau$ is an instance of $\sigma$'') holds when $\tau$ is obtained from
$\sigma$ by substituting order-0 types other than $\typ{o}$ for type variables. A \emph{first-order term} is a
$t$ with $\Sig;\Gamma\Vdash_f t:\tau$ for some $\tau$ of order~0 different from $\typ{o}$ \book{18}. The
first rule is where constants are polymorphic: the instance is chosen per occurrence. The second rule is
where variables are monomorphic: $\Gamma$ gives one type per variable.

The book states a principal typing property \book{18--19}: if a term has some typing, it has a most general
one, unique up to renaming of type variables. For example \lp{(X :: nil) :: nil} can be typed with
$X:\typ{int}$ or $X:\typ{list}\ \typ{int}$, and most generally with $X:A$, the whole term having type
\lp{list (list A)}. Application is left associative. In canonical form a term is a variable, a constant,
or $(f\ t_1\ \dots\ t_n)$ whose \emph{head} $f$ is always a constant \book{19}: first-order variables cannot
be applied.

\subsection{Representing data}\label{sec:horn-data}

The recipe for representing a class of objects is \book{19}: pick a type for it; declare one value
constructor per way of building an object, all with that common target type; let recursive components
have the target type as argument type. Binary trees \book{20}:
\begin{lprolog}
kind  btree  type -> type.
type  empty  btree A.
type  node   A -> btree A -> btree A -> btree A.
\end{lprolog}
\lp{(node 3 (node 1 empty empty) (node 4 empty empty))} has type \lp{btree int}, while a tree that mixes
\lp{3} and \lp{"cat"} is ill typed. This declaration is \emph{parametric}: every type variable of an
argument type also occurs in the target type. λProlog does not require this; the declaration
\lp{type node A -> btree -> btree -> btree.} gives heterogeneous trees \book{21}. Unlike ML datatypes, the
set of constructors of a type is \emph{open} (declarations can be added later), and terms can contain
variables, so a term describes a whole class of objects \book{21--22}.

The same recipe encodes the abstract syntax of a language. First-order formulas \book{22--23}:
\begin{lprolog}
kind  term, form  type.
type  a, b  term.
type  f     term -> term -> term.
type  var   string -> term.
% f(a, f(x, b))   is represented as   (f a (f (var "x") b))

type  p  term -> term -> form.
type  q  term -> form.

type  ff,        % encoding the false proposition
      tt         form.     % encoding the true proposition
type  &&,        % encoding conjunction
      !!,        % encoding disjunction
      ==>        form -> form -> form. % encoding implication
infixl  &&   5.
infixl  !!   4.
infixr  ==>  3.
\end{lprolog}
With the precedences above, \lp{(p a b) && (q a) !! (q (f a b)) ==> (p b (f b a))} stands for
$((p(a,b)\wedge q(a))\vee q(f(a,b)))\imp p(b,f(b,a))$.

The declarations \lp{infixl}, \lp{infixr} (and \lp{infix}, \lp{prefix}, \lp{postfix}, \dots) are
\emph{operator declarations} $\langle\text{fixity}\rangle\ c_1,\dots,c_n\ \langle\text{precedence}\rangle$:
they change only the concrete syntax, a higher number binding more tightly. For instance, after
\lp{infixl pr 5.} the expression \lp{(3 pr 4 pr "three")} stands for \lp{((pr ((pr 3) 4)) "three")}
\book{16}.

Imperative programs are encoded in the same way \book{24--25}. The book deliberately uses only two sorts,
\lp{expr} and \lp{stmt}, and one \emph{type-neutral} sort for all expressions, whether integer or boolean
valued; object-level typing is left to a later analysis written as a logic program:
\begin{lprolog}
kind stmt, expr  type.
type id     string -> expr.
type c      int -> expr.
type t, f   expr.
type &&, !!,
     plus, minus,
     mult, <    expr -> expr -> expr.
type :=     expr -> expr -> stmt.        % for assignment
type cond   expr -> stmt -> stmt -> stmt. % for conditionals
type while  expr -> stmt -> stmt.        % for while loops
type seq    stmt -> stmt -> stmt.        % for composition
infixl &&, mult   5.
infixl !!, plus, minus  4.
infix  <    3.
infix  :=   2.
\end{lprolog}
The C-like program \texttt{v = 1; i = n; while (0 < i) \{ v = v * i; i = i - 1; \}} is then the term
\begin{lprolog}
(seq ((id "v") := (c 1))
  (seq ((id "i") := (id "n"))
    (while ((c 0) < (id "i"))
      (seq ((id "v") := (id "v") mult (id "i"))
           ((id "i") := (id "i") minus (c 1))))))
\end{lprolog}
Because \lp{:=} has the lowest precedence, \lp{(id "v") := (id "v") mult (id "i")} parses as
\lp{(:= (id "v") (mult (id "v") (id "i")))}. The term \emph{is} the program's abstract syntax tree:
``first-order terms are a generalization of the labeled trees that are used in depicting abstract
syntax'' \book{25}. Types do not express every well-formedness condition (the left side of \lp{:=} must
encode an identifier); such conditions are checked by logic programs over the representation.

\begin{pitfall}[First-order syntax does not capture binding]
Adding \lp{type all term -> form -> form.} lets us write $\forall x\,(p(a,x)\wedge q(x))$ as
\lp{(all (var "x") ((p a (var "x")) && (q (var "x"))))}. The term has the right tree shape, but nothing
in it says that \lp{(var "x")} is \emph{bound} by \lp{all}: renaming, capture-avoiding substitution and
scoping must be programmed by hand in every computation over the syntax. ``This is a fundamental
limitation of first-order approaches to the treatment of syntax'' \book{23}. Representing object
variables by meta-level logic variables instead makes object scoping uncontrollable \book{22}. The book
removes the limitation with $\lambda$-tree syntax (\cref{ch:lambdatree}).
\end{pitfall}

\section{First-order unification}\label{sec:horn-unif}

Unification is the engine under backchaining. The book presents it as a nondeterministic rewriting of a
\emph{multiset of equations} \book{26--29}, the Martelli--Montanari style. A miniKanren programmer knows
the algorithm; the point of this section is the precise vocabulary (mgu, solved form, occurs check) and
the role of types.

\begin{remark}
The book is deliberately ``not a formal development of logic'' \book{6--7}: it states its results in
running prose, without numbered definitions or theorems. The definition, theorem and proposition
environments in this chapter are our packaging of those statements, each with the page where the book
makes it.
\end{remark}

Before the definitions, note the three ways a program uses unification \book{26--27}. Take
$X:\typ{int}$, $L:\typ{list}\ \typ{int}$.
\begin{itemize}
  \item \emph{Decomposition}: $\{(X\mathbin{:\!:}L)=(1\mathbin{:\!:}2\mathbin{:\!:}\cst{nil})\}$ has the unifier
    $\{\langle X,1\rangle,\langle L,2\mathbin{:\!:}\cst{nil}\rangle\}$: the left side is a \emph{pattern} that
    extracts head and tail.
  \item \emph{Recognition}: the same pattern unifies only with non-empty lists, and with $T$ repeated
    the pattern $(\cst{node}\ El\ T\ T)$ recognises the non-empty trees whose two subtrees are
    identical.
  \item \emph{Construction}: in
    \[\{(X\mathbin{:\!:}L_1)=(1\mathbin{:\!:}\cst{nil}),\ L_2=(2\mathbin{:\!:}\cst{nil}),\
    (X\mathbin{:\!:}L_2)=L_3\}\]
    the first equation takes an input apart, the last builds a new term from its
    pieces, and the ``output'' variable $L_3$ receives $1\mathbin{:\!:}2\mathbin{:\!:}\cst{nil}$.
\end{itemize}
The construction case is how a clause head such as \lp{append (X :: L1) L2 (X :: L3)} computes.

\begin{definition}[Unification problem, unifier, mgu {\book{26--28}}]\label{def:horn-mgu}
A \emph{unification problem} is a finite multiset of equations $t=s$ between first-order terms, the two
sides of each equation having the same type. A \emph{unifier} is a type-preserving substitution that
makes the two sides of every equation identical. A substitution $\theta$ is \emph{more general} than
$\varphi$ if $\varphi$ is obtained from $\theta$ by further substitution. A \emph{most general unifier} (mgu)
is a unifier more general than every other.
\end{definition}
Substitutions are \emph{constraints} on the values of variables. $\{\langle X,Y\rangle\}$ says only that
$X$ and $Y$ are equal; $\{\langle X,1\rangle,\langle Y,1\rangle\}$ satisfies that constraint but says more
\book{27}. A good unifier is as general as possible, because later steps of the computation will add
more constraints, and an over-specific choice could make them fail.

\begin{theorem}[mgu {\book{28}}]\label{thm:horn-mgu}
Typed first-order unification problems have most general unifiers whenever they have unifiers, and these
are unique up to substitutions that rename variables.
\end{theorem}
For instance the problem $\{(X\mathbin{:\!:}L)=(Y\mathbin{:\!:}Z\mathbin{:\!:}\cst{nil})\}$ has the two mgu's
\[
\{\langle X,Y\rangle,\langle L,Z\mathbin{:\!:}\cst{nil}\rangle\}
\quad\text{and}\quad
\{\langle Y,X\rangle,\langle L,Z\mathbin{:\!:}\cst{nil}\rangle\},
\]
which differ by the renaming $X\leftrightarrow Y$, and infinitely many less general ground unifiers \book{27--28}.

\paragraph{The transformation system.} The book gives three transformations, \emph{term reduction},
\emph{reorientation} and \emph{variable elimination} \book{28}. We split them into six rules
(this split is ours, the content is the book's). $E\uplus\{\dots\}$ is multiset union, $\bot$ is a problem
with no unifier, $\mathcal V$ the set of variables, and $f^\sigma$ means an occurrence of $f$ at type $\sigma$.
\begin{mathpar}
  \inferrule{E\uplus\{f^{\sigma}\,t_1\ldots t_n = f^{\sigma}\,s_1\ldots s_n\}}{E\uplus\{t_1=s_1,\ldots,t_n=s_n\}}\;\rulename{Decompose}
  \and
  \inferrule{E\uplus\{f^{\sigma}\,\vec t = g^{\sigma'}\,\vec s\}}{\bot}\;\rulename{Clash}\ (f\neq g\text{ or }\sigma\neq\sigma')
  \\
  \inferrule{E\uplus\{t=x\}}{E\uplus\{x=t\}}\;\rulename{Orient}\ (t\notin\mathcal V)
  \and
  \inferrule{E\uplus\{x=x\}}{E}\;\rulename{Delete}
  \and
  \inferrule{E\uplus\{x=t\}}{\bot}\;\rulename{Check}\ (x\in\operatorname{fv}(t),\ t\neq x)
  \\
  \inferrule{E\uplus\{x=t\}}{\subst{E}{t}{x}\uplus\{x=t\}}\;\rulename{Eliminate}\ (x\notin\operatorname{fv}(t),\ x\in\operatorname{fv}(E))
\end{mathpar}
\rulename{Decompose} and \rulename{Clash} are term reduction. They are justified because a \emph{rigid} term
$(f\ t_1\dots t_n)$ keeps its head under any substitution: different heads, \emph{or the same name at two
different types}, can never be made equal \book{28--29}. \rulename{Orient} is bookkeeping.
\rulename{Delete}, \rulename{Check} and \rulename{Eliminate} are variable elimination, and they rest on the
finiteness of terms: $x=t$ with $x$ strictly inside $t$ has no finite solution. This failure is the
\emph{occurs check}. The side condition of \rulename{Clash} is stated here for types without type
variables; with polymorphic constants, ``$\sigma\neq\sigma'$'' becomes ``$\sigma$ and $\sigma'$ have no common
instance'', and \rulename{Decompose} first unifies $\sigma$ with $\sigma'$ as generally as possible
\book{30--31}.

\begin{definition}[Solved form {\book{29}}]
$E$ is in \emph{solved form} if every equation has a variable on the left, and each such variable occurs
nowhere else in $E$. A solved form is read as a substitution, and it is its own mgu.
\end{definition}
\begin{theorem}[Termination and correctness {\book{29}}]
If variable elimination is never applied twice to the same unchanged equation, every sequence of
transformations terminates in $\bot$ or in a solved form, and every transformation preserves the set of
unifiers. The transformations thus define a nondeterministic algorithm for computing mgu's.
\end{theorem}

\begin{tutorial}[Three unification traces {\book{29--30}}]
\emph{Success.} The repeated variable $T$ makes the pattern recognise trees with two identical subtrees.
Write $n_2$ for $(\cst{node}\ 2\ \cst{empty}\ \cst{empty})$:
\begin{align*}
 & \{(\cst{node}\ El\ T\ T) = (\cst{node}\ 1\ n_2\ n_2)\}\\
 \longrightarrow\ & \{El=1,\ T=n_2,\ T=n_2\} && \text{term reduction}\\
 \longrightarrow\ & \{El=1,\ T=n_2,\ n_2=n_2\} && \text{variable elimination}\\
 \longrightarrow\ & \{El=1,\ T=n_2,\ 2=2,\ \cst{empty}=\cst{empty},\ \cst{empty}=\cst{empty}\} && \text{term reduction}\\
 \longrightarrow\ & \{El=1,\ T=n_2\} && \text{term reduction}\times 3
\end{align*}
The last step removes equations between constants: \rulename{Decompose} with $n=0$ arguments. The
result is a solved form, hence the mgu.

\emph{Constant clash.} If the right subtree is $(\cst{node}\ 3\ \cst{empty}\ \cst{empty})$ instead, the same
steps reach the equation $2=3$, and term reduction yields $\bot$.

\emph{Occurs check.} $\{T=(\cst{bt}\ 1\ T\ T)\}$ fails at once, because $T$ occurs in the right-hand side.
\end{tutorial}

\paragraph{Types in unification.} Types enter only in term reduction, when two heads have the same name:
their types must agree \book{30--31}. With polymorphic constants the type of an occurrence is not
determined by its name, so types cannot simply be checked once in advance. The book's example: with
$c:A\to A\to i$ and $a:i$, the problem $\{(c\ 1\ Y)=(c\ X\ a)\}$ is well typed (both sides have type $i$),
but in $(c\ 1\ Y)$ the constant $c$ occurs at type $\typ{int}\to\typ{int}\to i$ and in $(c\ X\ a)$ at
$i\to i\to i$. Ignoring types yields the \emph{ill-typed} unifier $\{\langle X,1\rangle,\langle Y,a\rangle\}$
\book{31}. The book stops there; completing the example, term reduction with types sees two occurrences
of $c$ at the types $\typ{int}\to\typ{int}\to i$ and $i\to i\to i$, which have no common instance, and
fails: the typed problem has no unifier. Types are
themselves first-order terms, so the same algorithm unifies them. The book concludes \book{31} that types
can be erased safely when the original problem's typing guarantees that the intermediate problems stay
well typed, in particular when every type variable in a constructor's argument types also appears in its
target type (the parametric condition of \cref{sec:horn-data}).

\paragraph{Logical reading and the occurs-check policy.} Constants behave like universally quantified
variables and logic variables like existentially quantified ones, so a unification problem is a formula
$\forall x_1\ldots\forall x_n\,\exists y_1\ldots\exists y_m\,[t_1=s_1\wedge\cdots\wedge t_p=s_p]$ \book{31--32}.
Later chapters mix the quantifier prefix (\cref{ch:hh,ch:hounif}). Many Prologs omit the occurs check for
speed; the book does not: ``we will require failure in any case where the occurs-check test is
satisfied'', because ``finiteness of terms, enforced through this test, will be important to encoding
quantifier dependencies'' \book{32}. Without it, deduction can be unsound \book{33}.

\begin{kanren}[Unification in \micro]
\micro's unifier already is the unifier of this section, minus types. Terms are Scheme data, so pairs,
vectors and records decompose structurally (\rulename{Decompose}), and different atoms or different shapes
fail (\rulename{Clash}). The substitution is \emph{triangular}: binding $x\mapsto t$ never rewrites the
other equations, and walking a term on demand plays the role of \rulename{Eliminate}. The occurs check is
on. Here are the book's problems, run on today's \micro{} (file \texttt{unify.scm}; each \scm{q} collects
the variables of interest, and each result is printed after its label):
\begin{scheme}
; U1  {(X :: L) = (1 :: 2 :: nil)}
(μkanren-run (q -1 #t) (fresh° (x l) (=° q (list x l)) (=° (cons x l) '(1 2))))
; U2  {(node El T T) = (node 1 (node 2 empty empty) (node 2 empty empty))}
(μkanren-run (q -1 #t)
  (fresh° (el t) (=° q (list el t))
    (=° `(node ,el ,t ,t) '(node 1 (node 2 empty empty) (node 2 empty empty)))))
; U3  the same with a 3 in the right subtree: constant clash
(μkanren-run (q -1 #t)
  (fresh° (el t) (=° q (list el t))
    (=° `(node ,el ,t ,t) '(node 1 (node 2 empty empty) (node 3 empty empty)))))
; U4  {(X :: L1) = (1 :: nil), L2 = (2 :: nil), (X :: L2) = L3}: construction
(μkanren-run (q -1 #t)
  (fresh° (x l1 l2) (=° (cons x l1) '(1)) (=° l2 '(2)) (=° (cons x l2) q)))
; U5  {(X :: L) = (Y :: Z :: nil)}: an mgu, shown up to renaming
(μkanren-run (q -1 #f)
  (fresh° (x l y z) (=° q (list x l y z)) (=° (cons x l) (list y z))))
; U6  {T = (bt 1 T T)}: occurs check
(μkanren-run (q -1 #t) (=° q `(bt 1 ,q ,q)))
; U7  {(c 1 Y) = (c X a)}: untyped unification returns the ill-typed unifier
(μkanren-run (q -1 #t) (fresh° (x y) (=° q (list x y)) (=° `(c 1 ,y) `(c ,x a))))
\end{scheme}
\begin{answers}
U1 ((1 (2)))
U2 ((1 (node 2 empty empty)))
U3 ()
U4 ((1 2))
U5 ((λ (α β) (cons α (cons (cons β (quote ())) (cons α (cons β (quote ())))))))
U6 ()
U7 ((1 a))
\end{answers}
In U4, \scm{q} plays the output variable $L_3$. U5 is run with \scm{#f} as the third argument of \scm{μkanren-run}, which prints the raw reified answer,
a \scm{λ} over the unbound variables, instead of evaluating it as \scm{#t} does; it reads $X=\alpha$, $L=(\beta)$, $Y=\alpha$, $Z=\beta$. Reification picks one
representative of the mgu ``up to renaming'' (\cref{thm:horn-mgu}). U7 is the warning of the book come
true: an untyped engine happily returns the ill-typed unifier. A typed layer would require typed
constant occurrences, or a static check plus erasure for parametric signatures. The proposal stays untyped
(\cref{sec:roadmap-lim-untyped}); simple types are a later work item (\cref{sec:roadmap-later-types}).
\end{kanren}

\section{First-order formulas}\label{sec:horn-formulas}

Book Chapter~2 relaxes the signature: constants may now have target type \lp{o}. They are the
\emph{predicate symbols}, and an expression of type \lp{o} headed by one of them is an \emph{atomic
formula} \book{34--35}:
\begin{lprolog}
type memb    A -> list A -> o.
type append  list A -> list A -> list A -> o.
\end{lprolog}
The logical constants and their λProlog spellings are \book{35}:
\begin{center}
\begin{tabular}{@{}lllll@{}}
\toprule
math & λProlog & type & fixity & conventional use\\
\midrule
$\top$ & \texttt{true} & \texttt{o} & & \\
$\imp$ & \texttt{=>} & \texttt{o -> o -> o} & infix, right assoc. & implication in goals\\
$\wedge$ & \texttt{\&} & \texttt{o -> o -> o} & infix, right assoc. & joins program clauses\\
$\wedge$ & \texttt{,} & \texttt{o -> o -> o} & infix, left assoc. & joins goals\\
$\vee$ & \texttt{;} & \texttt{o -> o -> o} & infix, left assoc. & \\
$\imp$ (reversed) & \texttt{:-} & \texttt{o -> o -> o} & infix, non assoc. & ``is implied by''\\
$\forall_\tau x\,F$ & \texttt{pi x\textbackslash\ F} & & binder & \\
$\exists_\tau x\,F$ & \texttt{sigma x\textbackslash\ F} & & binder & \\
\bottomrule
\end{tabular}
\end{center}
The two spellings of $\wedge$ and of $\imp$ are a stylistic convention \book{47--48}; breaking it does not
affect well-formedness. All these operators bind less tightly than application \book{35}, and \lp{&}
binds more tightly than \lp{:-} \book{48}, so \lp{A & B :- G} reads as $G\imp(A\wedge B)$. Only the
logical constants may take arguments of type \lp{o} in the first-order setting \book{35}. Quantifiers carry an optional type, \lp{(pi (x:T)\ F)}, inferred most generally when
omitted; the scope of \lp{x\ } extends as far to the right as possible, and only the binding pattern
matters, so \lp{pi x\ (p x) => sigma y\ (q x y, pi x\ (q y x))} and
\lp{pi z\ (p z) => sigma y\ (q z y, pi v\ (q y v))} are the same formula \book{36--37}. The typing rules
of the quantifiers (Figure~2.1, \book{38}) extend those of terms:
\begin{mathpar}
  \inferrule{\Sig;\Gamma,x:\tau\Vdash_f B:\typ{o}}{\Sig;\Gamma\Vdash_f \forall_\tau x\,B:\typ{o}}
  \and
  \inferrule{\Sig;\Gamma,x:\tau\Vdash_f B:\typ{o}}{\Sig;\Gamma\Vdash_f \exists_\tau x\,B:\typ{o}}
  \and
  \inferrule{\Sig;\Gamma\Vdash_f B:\tau}{\Sig;\Gamma\Vdash_f C:\tau}\ (B, C \text{ differ only in bound names})
\end{mathpar}
with the proviso, for the first two, that $x$ is not typed by $\Sig$ or $\Gamma$; the third rule renames
bound variables so that the proviso can be met. A \emph{first-order $\Sig$-formula} is an $F$ with
$\Sig;\emptyset\Vdash_f F:\typ{o}$; such formulas are closed \book{38}.

\section{The search semantics of goals}\label{sec:horn-search}

A logic programming language is fixed by four ingredients \book{38}: signatures, program clauses, goals, and
a calculus for building proofs. A sequent $\seq{\Sig}{\Prog}{G}$ has a set $\Prog$ of $\Sig$-formulas as its
program, and it states the judgment ``$G$ is provable from $\Prog$ and $\Sig$''.

\begin{definition}[Search semantics of the logical symbols {\book{39}}]\label{def:horn-reduction}
Read as steps of an interpreter:
\begin{description}[font=\normalfont\scshape]
  \item[and] reduce $\seq{\Sig}{\Prog}{B_1\wedge B_2}$ to both $\seq{\Sig}{\Prog}{B_1}$ and $\seq{\Sig}{\Prog}{B_2}$;
  \item[or] reduce $\seq{\Sig}{\Prog}{B_1\vee B_2}$ to either $\seq{\Sig}{\Prog}{B_1}$ or $\seq{\Sig}{\Prog}{B_2}$;
  \item[instan] reduce $\seq{\Sig}{\Prog}{\exists_\tau x\,B}$ to $\seq{\Sig}{\Prog}{\subst{B}{t}{x}}$ for some $\Sig$-term $t$ of type $\tau$;
  \item[augment] reduce $\seq{\Sig}{\Prog}{B_1\imp B_2}$ to $\seq{\Sig}{\Prog,B_1}{B_2}$;
  \item[generic] reduce $\seq{\Sig}{\Prog}{\forall_\tau x\,B}$ to $\seq{c:\tau,\Sig}{\Prog}{\subst{B}{c}{x}}$, where $c$ is a new constant, not in $\Sig$;
  \item[true] $\seq{\Sig}{\Prog}{\top}$ is provable at once.
\end{description}
\end{definition}
None of these looks at $\Prog$. Read downwards, they are the \emph{right-introduction rules} of a sequent
calculus, the book's Figure~2.2 \book{40}:
\begin{mathpar}
  \inferrule{ }{\seq{\Sig}{\Prog}{\top}}\;\top R
  \and
  \inferrule{\seq{\Sig}{\Prog}{B_1} \\ \seq{\Sig}{\Prog}{B_2}}{\seq{\Sig}{\Prog}{B_1\wedge B_2}}\;\wedge R
  \and
  \inferrule{\seq{\Sig}{\Prog}{B_1}}{\seq{\Sig}{\Prog}{B_1\vee B_2}}\;\vee R
  \and
  \inferrule{\seq{\Sig}{\Prog}{B_2}}{\seq{\Sig}{\Prog}{B_1\vee B_2}}\;\vee R
  \\
  \inferrule{\seq{\Sig}{\Prog,B_1}{B_2}}{\seq{\Sig}{\Prog}{B_1\imp B_2}}\;{\imp}R
  \and
  \inferrule{\seq{\Sig}{\Prog}{\subst{B}{t}{x}} \\ \Sig;\emptyset\Vdash_f t:\tau}{\seq{\Sig}{\Prog}{\exists_\tau x\,B}}\;\exists R
  \and
  \inferrule{\seq{c:\tau,\Sig}{\Prog}{\subst{B}{c}{x}}}{\seq{\Sig}{\Prog}{\forall_\tau x\,B}}\;\forall R
\end{mathpar}
($c$ not declared in $\Sig$ in $\forall R$). Each rule is sound in ``any interesting logic''. The question is
\emph{completeness} \book{41}: if a sequent with a non-atomic goal is provable, must the premise(s) of the
corresponding rule be provable too? The answer depends on which formulas are allowed and on the logic.

\begin{example}[Why goal-directedness can fail {\book{41--42}}]\label{ex:horn-counter}
Let $\Sig=\{p:\typ{o},q:\typ{o},r:i\to\typ{o},a:i,b:i\}$.
\begin{enumerate}
  \item $\seq{\Sig}{p\vee q}{q\vee p}$ is provable, but \rulename{or} leaves $\seq{\Sig}{p\vee q}{q}$ or
    $\seq{\Sig}{p\vee q}{p}$, and neither is provable. A disjunction in the \emph{program} breaks
    goal-directed search.
  \item $\seq{\Sig}{\emptyset}{p\vee(p\imp q)}$ is a classical tautology (excluded middle), but \rulename{or}
    leaves $p$, unprovable, or $p\imp q$, which \rulename{augment} reduces to the unprovable $\seq{\Sig}{p}{q}$.
  \item $\seq{\Sig}{(r\,a\wedge r\,b)\imp q}{\exists x\,(r\,x\imp q)}$ is classically provable by cases on
    $r\,a$, but every choice of witness $t$ leaves an unprovable sequent.
\end{enumerate}
\end{example}
The book concludes \book{42} that programs must not contain disjunctions (they must be \emph{definite},
``because they do not contain indefinite information''), and that the logic should be
\emph{intuitionistic}, where excluded middle does not hold. Counterexamples~2 and~3 need $\imp$ in the
\emph{goal}, which Horn clauses forbid; the switch of logic pays off only from \cref{ch:hh} on.

Atomic goals need rules that look at the program; the general framework leaves them open on purpose
\book{42}. When the search succeeds, the system reports an \emph{answer substitution}: the free variables
of a query are read as existentially quantified at its head, and the answer maps them to terms that make
the computation succeed \book{42--43}.

\section{Horn clauses and backchaining}\label{sec:horn-fohc}

\begin{definition}[fohc {\book{43}}]\label{def:horn-fohc}
With $A$ ranging over atomic formulas, the goals $G$ and program clauses $D$ of \emph{first-order Horn
clauses} (\fohc) are
\[
  G ::= \top \mid A \mid G\wedge G \mid G\vee G \mid \exists_\tau x\,G
  \qquad
  D ::= A \mid G\imp D \mid D\wedge D \mid \forall_\tau x\,D .
\]
\end{definition}
Only $\imp$ and $\forall$ are missing at the top of goals, and only $\vee$ and $\exists$ at the top of
clauses. The family contains the familiar Prolog shape $\forall x_1\ldots\forall x_m\,(A_1\wedge\cdots\wedge
A_n\imp A_0)$, but is more liberal: the head may sit under further implications, conjunctions and
universals \book{43--44}.

Backchaining is first explained informally \book{44}: if $\Prog$ contains $A$, succeed; if it contains $G'\imp
A$, reduce the goal $A$ to $G'$. Formally it uses a \emph{focused} sequent $\Sig;\Prog\xrightarrow{D}A$, ``try
to prove the atom $A$ by backchaining on $D$''. The rules of the book's Figure~2.3 \book{45} are
\begin{mathpar}
  \inferrule{\Sig;\Prog\xrightarrow{D}A}{\seq{\Sig}{\Prog}{A}}\;\rulename{decide}\ (D\in\Prog)
  \and
  \inferrule{ }{\Sig;\Prog\xrightarrow{A}A}\;\rulename{initial}
  \and
  \inferrule{\Sig;\Prog\xrightarrow{D}A \\ \seq{\Sig}{\Prog}{G}}{\Sig;\Prog\xrightarrow{G\imp D}A}\;{\imp}L
  \\
  \inferrule{\Sig;\Prog\xrightarrow{D_1}A}{\Sig;\Prog\xrightarrow{D_1\wedge D_2}A}\;\wedge L
  \and
  \inferrule{\Sig;\Prog\xrightarrow{D_2}A}{\Sig;\Prog\xrightarrow{D_1\wedge D_2}A}\;\wedge L
  \and
  \inferrule{\Sig;\Prog\xrightarrow{\subst{D}{t}{x}}A \\ \Sig;\emptyset\Vdash_f t:\tau}{\Sig;\Prog\xrightarrow{\forall_\tau x\,D}A}\;\forall L
\end{mathpar}
\rulename{decide} picks a clause; \rulename{initial} closes the branch when the focused atom is
\emph{identical} to the goal; ${\imp}L$ proves the body $G$ and continues on $D$; $\wedge L$ picks a
conjunct; $\forall L$ instantiates. The formula over the arrow is on the left of the sequent, in focus, so
these are \emph{left}-introduction rules complementing Figure~2.2. For a Prolog-shaped clause
$D=\forall\bar x\,(A_1\wedge\cdots\wedge A_n\imp A_0)$ the rules compose into one step \book{45}:
\begin{mathpar}
  \inferrule{\seq{\Sig}{\Prog}{A_1\theta} \\ \cdots \\ \seq{\Sig}{\Prog}{A_n\theta}}{\Sig;\Prog\xrightarrow{D}A}
\end{mathpar}
where $\theta$ maps each $x_i$ to a $\Sig$-term of its type and $A=A_0\theta$: the instance
$A_1\theta\wedge\cdots\wedge A_n\theta\imp A_0\theta$ follows from $\Prog$, and modus ponens gives $A_0\theta$.
This grouping of small rules into large ones is a first instance of Andreoli's \emph{focusing}
\cite{Andreoli1992}, which the book cites as the general setting \book{73}.

\begin{definition}[O-proof {\book{45}}]
An \emph{O-proof} is a proof built with the rules of the book's Figures~1.2, 2.2 and~2.3.
\end{definition}
\begin{theorem}[Completeness of goal-directed search for fohc {\book{45}}]\label{thm:horn-oproof}
If $\Prog$ is an \fohc{} program and $G$ an \fohc{} goal, then $\seq{\Sig}{\Prog}{G}$ has an O-proof iff it is
provable in classical logic iff it is provable in intuitionistic logic.
\end{theorem}
\begin{proposition}[Flatness {\book{46}}]\label{prop:horn-flat}
Every sequent in an O-proof of $\seq{\Sig}{\Prog}{G}$ has the form $\seq{\Sig}{\Prog}{G'}$ or
$\Sig;\Prog\xrightarrow{D}A$, with the \emph{same} $\Sig$ and $\Prog$.
\end{proposition}
So in \fohc{} the program and signature are flat and global, and \rulename{augment} and \rulename{generic}
never fire. These two rules are exactly what \cref{ch:hh} adds. The book's term for a logic in which this
kind of goal-directed (``uniform'') search is complete is \emph{abstract logic programming
language}~\cite{MillerNadathurPfenningScedrov1991,MillerNadathur2012} \book{73}; \fohc{} is the base case.

\begin{tutorial}[An O-proof for a query with two answers]
Take the program $\Prog=\{D_1,D_2\}$ of \lp{append} (\cref{sec:horn-syntax}), with
$D_1=\forall l\,(\cst{app}\ \cst{nil}\ l\ l)$ and $D_2=\forall x\forall l_1\forall l_2\forall l_3\,(\cst{app}\ l_1\
l_2\ l_3\imp\cst{app}\ (x\mathbin{:\!:}l_1)\ l_2\ (x\mathbin{:\!:}l_3))$, and the goal $\exists X\exists Y\,\cst{app}\ X\ Y\
(1\mathbin{:\!:}\cst{nil})$. Instead of guessing terms in $\exists R$ and $\forall L$ we use logic variables and let
\rulename{initial} unify (this is the practical version of the rules, \cref{sec:horn-pragmatics}).
\begin{enumerate}
  \item $\exists R$ twice: logic variables $X,Y$; goal $\cst{app}\ X\ Y\ (1\mathbin{:\!:}\cst{nil})$, atomic.
  \item \rulename{decide} $D_1$; $\forall L$ with a new $L$; \rulename{initial} unifies
    $\cst{app}\ \cst{nil}\ L\ L$ with the goal: $X=\cst{nil}$, $L=Y$, $Y=1\mathbin{:\!:}\cst{nil}$.
    \emph{First answer}: $X=\cst{nil}$, $Y=1\mathbin{:\!:}\cst{nil}$.
  \item Back to step 2, \rulename{decide} $D_2$ instead. $\forall L$ four times with new $X',L_1,L_2,L_3$,
    then ${\imp}L$ splits into a focused premise and a goal premise. \rulename{initial} on the focused
    premise unifies $\cst{app}\ (X'\mathbin{:\!:}L_1)\ L_2\ (X'\mathbin{:\!:}L_3)$ with the goal: $X=X'\mathbin{:\!:}L_1$, $Y=L_2$,
    $X'=1$, $L_3=\cst{nil}$. The goal premise is now $\cst{app}\ L_1\ Y\ \cst{nil}$.
  \item For $\cst{app}\ L_1\ Y\ \cst{nil}$: \rulename{decide} $D_1$ succeeds with $L_1=\cst{nil}$,
    $Y=\cst{nil}$. \emph{Second answer}: $X=1\mathbin{:\!:}\cst{nil}$, $Y=\cst{nil}$. \rulename{decide} $D_2$ fails:
    its head needs the third argument to be $x\mathbin{:\!:}l_3$, which clashes with $\cst{nil}$.
\end{enumerate}
The O-proof for the second answer, with the unifier already applied. Write $A_0$ for
$\cst{app}\ (1\mathbin{:\!:}\cst{nil})\ \cst{nil}\ (1\mathbin{:\!:}\cst{nil})$ and $A_1$ for $\cst{app}\ \cst{nil}\ \cst{nil}\ \cst{nil}$:
\begin{mathpar}
\inferrule*[right=$\exists R\times 2$]{
 \inferrule*[right=decide]{
  \inferrule*[right=$\forall L\times 4$]{
   \inferrule*[right=${\imp}L$]{
     \inferrule*[right=initial]{ }{\Sig;\Prog\xrightarrow{A_0}A_0}
     \\
     \inferrule*[right=decide$+\forall L+$initial]{ }{\seq{\Sig}{\Prog}{A_1}}
   }{\Sig;\Prog\xrightarrow{A_1\imp A_0}A_0}
  }{\Sig;\Prog\xrightarrow{D_2}A_0}
 }{\seq{\Sig}{\Prog}{A_0}}
}{\seq{\Sig}{\Prog}{\exists X\exists Y\,\cst{app}\ X\ Y\ (1\mathbin{:\!:}\cst{nil})}}
\end{mathpar}
Every sequent in it has the same $\Sig$ and $\Prog$, as \cref{prop:horn-flat} promises.
\end{tutorial}

\section{λProlog concrete syntax and interaction}\label{sec:horn-syntax}

In program clauses, an unbound token beginning with an uppercase letter is a variable universally
quantified over the whole clause; \lp{_} is an anonymous variable, each occurrence distinct \book{46--47}.
So \book{47}
\begin{lprolog}
type append  list A -> list A -> list A -> o.
append nil L L.
append (X :: L1) L2 (X :: L3) :- append L1 L2 L3.
\end{lprolog}
abbreviates \lp{pi L\ append nil L L.} and
\lp{pi X\ pi L1\ pi L2\ pi L3\ append (X :: L1) L2 (X :: L3) :- append L1 L2 L3.}; the types of the
variables are inferred (\lp{X : A}, the lists \lp{list A}). Anonymous variables save names for values
that are never used again \book{47}:
\begin{lprolog}
type sublist  list A -> list A -> o.
sublist L K :- append _ T K, append L _ T.
\end{lprolog}
is the same clause as \lp{sublist L K :- append U T K, append L V T.}: $L$ is a sublist of $K$ if $K$
has a suffix $T$ of which $L$ is a prefix.

\paragraph{Equivalent groupings of clauses.} Because the D-grammar allows $\wedge$ and $\imp$ to nest, one
predicate can be written in several equivalent ways \book{48--49}. With
\lp{kind bool type.}, constructors \lp{neg}, \lp{and}, \lp{or}, \lp{imp} and \lp{type ident bool -> bool -> o.},
the four clauses of Form~1 can be compressed into a single clause, Form~3:
\begin{lprolog}
ident (neg B) (neg D) :- ident B D.
ident (and B C) (and D E) :- ident B D, ident C E.
ident (or  B C) (or  D E) :- ident B D, ident C E.
ident (imp B C) (imp D E) :- ident B D, ident C E.
\end{lprolog}
\begin{lprolog}
ident (neg B) (neg D) &
 (ident (and B C) (and D E) &
  ident (or  B C) (or  D E) &
  ident (imp B C) (imp D E) :- ident C E) :- ident B D.
\end{lprolog}
How does Form~3 parse and run? Since \lp{&} binds more tightly than \lp{:-}, it is the single D-formula
\[
  G_1\imp\bigl(D_{\cst{neg}}\wedge(G_2\imp(D_{\cst{and}}\wedge D_{\cst{or}}\wedge D_{\cst{imp}}))\bigr)
\]
(universally closed), where $G_1$ stands for the goal \lp{ident B D}, $G_2$ for
\lp{ident C E}, and $D_{\cst{neg}},\dots$ for the four heads. It is legal because the D-grammar allows
$G\imp D$ and $D\wedge D$ to nest. To solve an atom \lp{ident (and B0 C0) X}, \rulename{decide} picks this
clause, $\forall L$ instantiates its variables, ${\imp}L$ sets the body \lp{ident B D} aside as a goal
premise and keeps backchaining on the conjunction, $\wedge L$ chooses the right conjunct, ${\imp}L$ sets
\lp{ident C E} aside, $\wedge L$ chooses $D_{\cst{and}}$, and \rulename{initial} unifies it with the goal.
The two goals set aside are exactly the body of the second clause of Form~1. The book claims that the
forms are equivalent in derivable goals and in observable behaviour, and that a smart implementation
preprocesses Form~3 into Form~1 \book{49}. (Figure~2.3 does not say in which order the two premises of
${\imp}L$ are explored; for the behaviour to match Form~1, whose body is \lp{ident B D, ident C E}, the
outer premise must be solved first. This remark is ours.) (As printed, \lp{ident} has no base case,
so it proves no closed goal; the example is about clause structure only.)

\paragraph{The read--prove--print loop.} Programs live in modules (\cref{sec:hoprog-modules}). Teyjus compiles a
module with \texttt{tjcc}, links it with \texttt{tjlink} and runs it with \texttt{tjsim}, which shows the
prompt \lp{[lists] ?-} \book{281--282}. In \emph{queries} the convention for variables is dual: an
uppercase token is \emph{existentially} quantified, with outermost scope \book{50}. The prove phase looks
for an O-proof, with three possible outcomes: a proof, no proof, or non-termination. Non-termination cannot
be ruled out by any interpreter, because provability in \fohc{} is undecidable \book{51}. The print phase
shows the bindings of the named query variables only \book{51--53}:
\begin{lprolog}
[lists] ?- append (1::nil) (2::nil) (3::nil).
no
[lists] ?- append (1::nil) (2::nil) (1::2::nil).
solved
[lists] ?- append (1::nil) (2::nil) X.
X = (1::2::nil)
[lists] ?- append X Y (1::nil).
X = nil
Y = 1::nil;
X = 1::nil
Y = nil;
no
\end{lprolog}
Typing \lp{;} asks for another proof. The queries \lp{sigma X\ append (1::nil) (2::nil) X} and
\lp{append (1::nil) (2::nil) _} attempt the same sequent as the third query, but print just
\lp{solved}: \lp{sigma}-bound and anonymous variables are not shown \book{51--52}.

\begin{kanren}[The same session in \micro]
A \fohc{} predicate becomes a \scm{define-relation} whose body is a \scm{cond°} with one line per clause:
\scm{fresh°} for the clause variables, \scm{=°} for head unification, then the body goals. File
\texttt{append.scm}:
\begin{scheme}
;; append nil L L.
;; append (X :: L1) L2 (X :: L3) :- append L1 L2 L3.
(define-relation (append° l k m)
  (cond° ((=° l '()) (=° k m))
         ((fresh° (x l1 l3)
            (=° l (cons x l1)) (=° m (cons x l3))
            (append° l1 k l3)))))

(μkanren-run (q -1 #t) (append° '(1) '(2) '(3)))
(μkanren-run (q -1 #t) (append° '(1) '(2) '(1 2)))
(μkanren-run (q -1 #t) (append° '(1) '(2) q))
(μkanren-run (q -1 #t) (fresh° (x y) (=° q (list x y)) (append° x y '(1))))
\end{scheme}
\begin{answers}
append (1) (2) (3)    => ()
append (1) (2) (1 2)  => (α)
append (1) (2) X      => ((1 2))
append X Y (1)        => ((() (1)) ((1) ()))
\end{answers}
\scm{()} is \lp{no}. The answer \scm{(α)} is \lp{solved}: one proof, with the query variable left
unconstrained. The run argument \scm{-1} asks for all answers, what typing \lp{;} repeatedly does in
Teyjus. The variables of an inner \scm{fresh°} are hidden like \lp{sigma}-bound ones.
\end{kanren}

\section{Core examples}\label{sec:horn-examples}

The book illustrates three programming styles \book{53--59}: search problems, relations over recursive
data, and relations over abstract syntax.

\subsection{Search: finite-state machines}\label{sec:horn-fsm}

The book's Figure~2.4 \book{54} describes acceptance by a finite-state machine whose arcs are labelled by
\emph{words}:
\begin{lprolog}
kind state, letter        type.
type q1, q2, q3, q4, q5   state.
type a, b                 letter.
type start, final         state -> o.
type path, trans          state -> list letter -> state -> o.
type accept               list letter -> o.

path S nil S.
path S Letters T :- trans S Arc M, append Arc Rest Letters,
                    path M Rest T.
accept W :- start S, path S W F, final F.
\end{lprolog}
The goal \lp{append Arc Rest Letters} consumes the arc's word as a prefix of the input. Figure~2.5 is a
nondeterministic machine for $(a+b)^*(ab+ba)$, the words ending in $ab$ or $ba$:
\begin{lprolog}
start q1 & final q2 & final q3.
trans q1 (a::nil) q1      & trans q1 (b::nil) q1.
trans q1 (a::b::nil) q2   & trans q1 (b::a::nil) q3.
\end{lprolog}
To enumerate accepted words, the book first generates lists of increasing length \book{55}:
\begin{lprolog}
type lists  list A -> o.
lists nil.
lists (_::L) :- lists L.
\end{lprolog}
and reports \book{55--56}
\begin{lprolog}
[fsm1] ?- accept (b::b::a::b::nil).
solved
[fsm1] ?- accept (b::a::X::Y::nil).
X = a
Y = b ;
X = b
Y = a ;
no
[fsm1] ?- lists L.
L = nil ;
L = T1 :: nil ;
L = T1 :: T2 :: nil ;
L = T1 :: T2 :: T3 :: nil
[fsm1] ?- lists L, accept L.
L = a :: b :: nil ;
L = b :: a :: nil ;
L = a :: a :: b :: nil ;
L = a :: b :: a :: nil ;
L = b :: a :: b :: nil
\end{lprolog}
Of the last two queries the book notes that ``both these queries have an unbounded number of successful
proofs''; a depth-first engine ``will systematically produce lists of increasing length'' \book{55--56}.
Note the answer \lp{L = T1 :: nil}: an answer substitution may contain variables, and then every instance
of it is an answer. Why is the generator needed? Under
depth-first search, \lp{accept W} with \lp{W} unbound tries the first \lp{trans} clause
\lp{trans q1 (a::nil) q1} first; \lp{append} then binds \lp{W = a :: Rest}, and \lp{path q1 Rest F} is
in the same situation again, forever. No answer is ever printed. \lp{lists L} breaks the loop by giving
\lp{accept} lists of fixed length, one length at a time.

\begin{pitfall}[An erratum in Figure~2.6]
The book also gives a deterministic five-state machine (Figure~2.6, \book{54}), printed as
\begin{lprolog}
start q1 & final q4 & final q5.
trans q1 (a::nil) q2   &   trans q1 (b::nil) q3.
trans q2 (a::nil) q1   &   trans q2 (b::nil) q4.
trans q3 (a::nil) q5   &   trans q3 (b::nil) q1.
trans q4 (a::nil) q5   &   trans q4 (b::nil) q3.
trans q5 (a::nil) q2   &   trans q5 (b::nil) q4.
\end{lprolog}
and says that exactly the same behaviour would be observed with it \book{56}. As printed it does not
accept the stated language: the
clauses \lp{trans q2 (a::nil) q1} and \lp{trans q3 (b::nil) q1} send the machine back to the start, so it
rejects \lp{aab}, \lp{bba}, and so on. Replacing them by \lp{trans q2 (a::nil) q2} and
\lp{trans q3 (b::nil) q3} fixes it. A brute-force comparison with the regular expression on all words of
length up to~8 confirms both the defect and the fix.
\end{pitfall}

\subsection{Recursive data: insertion into a search tree}

Relations over recursive data follow the structure of the data \book{57}:
\begin{lprolog}
type insert  int -> btree int -> btree int -> o.
insert X empty (node X empty empty).
insert X (node A L R) (node A NL R) :- X < A,  insert X L NL.
insert X (node A L R) (node A L NR) :- X >= A, insert X R NR.
\end{lprolog}
\begin{lprolog}
[btree] ?- insert 4 (node 3 (node 2 empty empty) empty) T.
T = node 3 (node 2 empty empty) (node 4 empty empty)
\end{lprolog}
The comparisons \lp{<} and \lp{>=} are built-ins ``correctly used only when both arguments are
instantiated to numbers'' \book{49}, so \lp{insert} is well defined only when its first two arguments are
closed terms \book{57}. This is a \emph{mode} restriction: the relation is not usable in every direction.

\subsection{Abstract syntax: a sequent prover}\label{sec:horn-prv}

The third style uses \fohc{} to animate an inference system. The book encodes a classical propositional
sequent calculus with multisets $\Gamma\longrightarrow\Delta$ on both sides (Figure~2.7, \book{58}), over the
formulas of \cref{sec:horn-data}:
\begin{mathpar}
  \inferrule{ }{\Gamma,\bot\longrightarrow\Delta}\;\bot L
  \and
  \inferrule{ }{\Gamma,A\longrightarrow\Delta,A}\;\rulename{initial}
  \\
  \inferrule{\Gamma,A\longrightarrow\Delta \\ \Gamma,B\longrightarrow\Delta}{\Gamma,A\vee B\longrightarrow\Delta}\;\vee L
  \and
  \inferrule{\Gamma\longrightarrow\Delta,A,B}{\Gamma\longrightarrow\Delta,A\vee B}\;\vee R
  \\
  \inferrule{\Gamma,A,B\longrightarrow\Delta}{\Gamma,A\wedge B\longrightarrow\Delta}\;\wedge L
  \and
  \inferrule{\Gamma\longrightarrow\Delta,A \\ \Gamma\longrightarrow\Delta,B}{\Gamma\longrightarrow\Delta,A\wedge B}\;\wedge R
  \\
  \inferrule{\Gamma\longrightarrow\Delta,A \\ \Gamma,B\longrightarrow\Delta}{\Gamma,A\imp B\longrightarrow\Delta}\;{\imp}L
  \and
  \inferrule{\Gamma,A\longrightarrow\Delta,B}{\Gamma\longrightarrow\Delta,A\imp B}\;{\imp}R
\end{mathpar}
Several formulas on the right make the calculus classical. Here $\Gamma,\Delta,A,B$ are \emph{schema}
variables of the object calculus, not logic variables of λProlog. Multisets are lists, and a rule matches
``$\Gamma, P$'' by selecting $P$ and keeping the rest:
\begin{lprolog}
type memb_and_rest  A -> list A -> list A -> o.
memb_and_rest X (X :: L) L.
memb_and_rest X (Y :: L) (Y :: L1) :- memb_and_rest X L L1.
\end{lprolog}
The translation recipe is mechanical \book{59}: \emph{a rule's conclusion becomes the head of a clause, its
premises become the body, and the schema variables become clause variables}. The eight clauses, one per
rule in the order $\bot L$, \rulename{initial}, $\wedge L$, $\vee L$, ${\imp}L$, $\wedge R$, $\vee R$,
${\imp}R$, as printed:
\begin{lprolog}
type prv  list form -> list form -> o.

prv Gamma Delta :- memb_and_rest ff Gamma _.
prv Gamma Delta :- memb_and_rest A Gamma _,
                   memb_and_rest A Delta _.
prv Gamma Delta :-
      memb_and_rest (A && B) Gamma Gamma',
      prv (A :: B :: Gamma') Delta.
prv Gamma Delta :-
      memb_and_rest (A !! B) Gamma Gamma',
      prv (A :: Gamma') Delta, prv (B :: Gamma') Delta.
prv Gamma Delta :-
      memb_and_rest (A ==> B) Gamma Gamma',
      prv Gamma' (A :: Delta), prv (B :: Gamma') Delta.
prv Gamma Delta :-
      memb_and_rest (A && B) Delta Delta',
      prv Gamma (A :: Delta'), prv Gamma (B :: Delta').
prv Gamma Delta :-
      memb_and_rest (A !! B) Delta Delta',
      prv Gamma (A :: B :: Delta').
prv Gamma Delta :-
      memb_and_rest (A ==> B) Delta Delta',
      prv (A :: Delta) (B :: Gamma').
\end{lprolog}
Read the $\vee L$ clause against its rule: the conclusion $\Gamma,A\vee B\longrightarrow\Delta$ becomes the
head \lp{prv Gamma Delta} together with the selection \lp{memb_and_rest (A !! B) Gamma Gamma'} (so the
clause's \lp{Gamma'} is the rule's $\Gamma$), and the two premises become the two body goals. The
\rulename{initial} clause needs no selection of the rest: it only asks that some $A$ be on both sides.
\begin{lprolog}
[logic] ?- prv nil (((p a b) !! ((p a b) ==> (q a a))) :: nil).
solved
\end{lprolog}
Note the \emph{level shift}. This is $p\vee(p\imp q)$, the formula that \cref{ex:horn-counter} showed
goal-directed search cannot prove (it is not even an \fohc{} goal, since it contains $\imp$). Here it is
object-level \emph{data}, proved by an object-level classical calculus that is \emph{programmed} in
\fohc: the search proves the meta-level atom \lp{prv nil (...)}, by $\vee R$ and then ${\imp}R$ and
\rulename{initial} on \lp{p a b}.

\begin{pitfall}[An erratum in the last \texttt{prv} clause]
The clause for ${\imp}R$ is printed \book{59} with the body \lp{prv (A :: Delta) (B :: Gamma')}. It should be
\lp{prv (A :: Gamma) (B :: Delta')}. As printed, \lp{Gamma'} is a singleton variable, the right-hand side
becomes an open list, and the initial clause can unify a hypothesis into that open tail. The printed
program then ``proves'' $p\,a\,b\imp q\,a\,a$, which is not valid. A \micro{} port of the printed clause
confirms it: with the printed clause, the query below for \lp{(p a b) ==> (q a a)} returns \scm{(α)}
instead of \scm{()} (file \texttt{prv-printed.scm}).
\end{pitfall}

\begin{kanren}[The prover in \micro]
The port keeps the clause-per-line shape (file \texttt{prv.scm}, last clause corrected). Formulas are
lists such as \scm{(!! A B)}:
\begin{scheme}
(define-relation (memb-and-rest° x l r)
  (cond° ((=° l (cons x r)))
         ((fresh° (y l1 r1)
            (=° l (cons y l1)) (=° r (cons y r1)) (memb-and-rest° x l1 r1)))))

(define-relation (prv° Γ Δ)
  (cond°
    ((fresh° (_) (memb-and-rest° 'ff Γ _)))                                   ; ⊥L
    ((fresh° (a _1 _2) (memb-and-rest° a Γ _1) (memb-and-rest° a Δ _2)))      ; initial
    ((fresh° (a b Γ*) (memb-and-rest° `(&& ,a ,b) Γ Γ*)                       ; ∧L
       (prv° `(,a ,b . ,Γ*) Δ)))
    ((fresh° (a b Γ*) (memb-and-rest° `(!! ,a ,b) Γ Γ*)                       ; ∨L
       (prv° (cons a Γ*) Δ) (prv° (cons b Γ*) Δ)))
    ((fresh° (a b Γ*) (memb-and-rest° `(==> ,a ,b) Γ Γ*)                      ; ⊃L
       (prv° Γ* (cons a Δ)) (prv° (cons b Γ*) Δ)))
    ((fresh° (a b Δ*) (memb-and-rest° `(&& ,a ,b) Δ Δ*)                       ; ∧R
       (prv° Γ (cons a Δ*)) (prv° Γ (cons b Δ*))))
    ((fresh° (a b Δ*) (memb-and-rest° `(!! ,a ,b) Δ Δ*)                       ; ∨R
       (prv° Γ `(,a ,b . ,Δ*))))
    ((fresh° (a b Δ*) (memb-and-rest° `(==> ,a ,b) Δ Δ*)                      ; ⊃R
       (prv° (cons a Γ) (cons b Δ*))))))

(define pab '(p a b))
(define qaa '(q a a))
(μkanren-run (q 1 #t) (prv° '() `((!! ,pab (==> ,pab ,qaa)))))
(μkanren-run (q 1 #t) (prv° '() `((==> ,pab ,qaa))))
\end{scheme}
\begin{answers}
prv nil (p a b !! (p a b ==> q a a))   => (α)
prv nil (p a b ==> q a a)              => ()
\end{answers}
\end{kanren}

\section{Pragmatics of the search}\label{sec:horn-pragmatics}

The rules of \cref{sec:horn-search,sec:horn-fohc} are nondeterministic. λProlog resolves the choices with
a ``simple and rigid search strategy'' so that a programmer can predict behaviour and resource usage
\book{60}:
\begin{enumerate}
  \item for $\vee R$ and $\wedge L$, the left subformula is tried before the right one;
  \item for $\wedge R$, the left premise is derived first;
  \item for \rulename{decide}, $\Prog$ is a \emph{list}, not a set: clauses are tried in textual order, and
    order and multiplicity matter;
  \item the terms of $\exists R$ and $\forall L$ are not enumerated (most types are infinite) but replaced
    by \emph{logic variables}, instantiated by unification \book{61}.
\end{enumerate}
The price is \emph{incompleteness}: ``there may be sequents that have derivations but for which no
derivations will be found because of the strategy. In fact, the chosen strategy even can cause indefinite
looping when a more flexible search strategy might be able to find a proof'' \book{60}. Since the strategy
is known, programs can usually be restructured, as \lp{lists L} did for \lp{accept}.

With a logic variable $X$ that is new in the search, the quantifier rules and \rulename{initial} become
\book{61}
\begin{mathpar}
  \inferrule{\Sig;\emptyset\Vdash_f X:\tau \\ \seq{\Sig}{\Prog}{\subst{B}{X}{x}}}{\seq{\Sig}{\Prog}{\exists_\tau x\,B}}
  \and
  \inferrule{\Sig;\Prog\xrightarrow{\subst{D}{X}{x}}A \\ \Sig;\emptyset\Vdash_f X:\tau}{\Sig;\Prog\xrightarrow{\forall_\tau x\,D}A}
  \and
  \inferrule{ }{\Sig;\Prog\xrightarrow{A'}A}\;\rulename{initial}\ (A, A'\text{ unifiable})
\end{mathpar}
where the unifier must be ``percolated all over the derivation that has been constructed thus far'', and
typed unification guarantees that the delayed typing premises hold. Finally the program can be precompiled
\book{61--62}: applying ${\imp}L$ and $\wedge L$ in advance turns it into a list of clauses
$\forall\bar x\,(G_1\wedge\cdots\wedge G_n\imp A)$, which are then partitioned by predicate name, compiled,
and indexed on their arguments, as in the Warren Abstract Machine.

\paragraph{Cut, negation as failure, arithmetic.} Book Chapter~2 has no pruning operator; it only
stresses that the cut \emph{rule} ``has no connection to Prolog's pruning operator !'' \book{63}. Teyjus
provides \lp{!}, \lp{fail} and \lp{not} as metalogical predicates, with the warning that
``negation-by-failure has a problematic semantics in a situation where implicational goals are used in an
essential way'' \book{287}. The book later defines \book{137}
\begin{lprolog}
type not  o -> o.
not P :- P, !, fail.
not P.
\end{lprolog}
and shows how far this is from logic: \lp{?- X = 2, not (1 = X).} succeeds, whereas
\lp{?- not (1 = X), X = 2.} fails, so ``even a basic property such as commutativity of conjunctions becomes
invalid'' \book{137--138} (\cref{ch:hoprog}). Arithmetic is \emph{intensional}: \lp{2 + 3} is a term, so the goal
\lp{(2 + 3 = 5)} fails \book{66}, and only the non-logical predicate \lp{is} evaluates
\book{67,286}. The comparisons used by \lp{insert} are of the same non-logical kind.

\section{Logical notions: cut and presentations of fohc}\label{sec:horn-logic}

\Cref{thm:horn-oproof} gives \fohc{} the duality the book wants: solving a goal means both ``it follows
logically'' and ``a certain derivation exists'' \book{62}. The cut rule \book{63} is
\begin{mathpar}
  \inferrule{\seq{\Sig}{\Prog}{B} \\ \seq{\Sig}{\Prog,B}{G}}{\seq{\Sig}{\Prog}{G}}\;\rulename{cut}
\end{mathpar}
It uses a lemma $B$. Choosing lemmas takes ingenuity, so cut is not used to compute, but it is used to
\emph{reason about} computation. Cut elimination says that adding cut to classical or intuitionistic
logic does not change the derivable sequents, and the completeness of O-proofs means that an adapted cut is
admissible for them. Consequence: replacing a subformula of a program by a logically equivalent one gives a
program that proves the same goals, ``in the sense of O-provability, although maybe not under a
depth-first search strategy'' \book{63}.

This licenses three equivalent presentations of \fohc{} clauses \book{63--64}: the liberal one of
\cref{def:horn-fohc}; the common Prolog one, $F ::= A\mid F\wedge F$, $D ::= A\mid F\imp A\mid\forall_\tau
x\,D$; and a compact one, $D ::= A\mid A\imp D\mid\forall_\tau x\,D$. A clause of one presentation is
equivalent to a set of clauses of another through equivalences such as
\[
\forall x\,(B_1\wedge B_2)\equiv(\forall x\,B_1)\wedge(\forall x\,B_2),\qquad
B_1\imp(B_2\imp B_3)\equiv(B_1\wedge B_2)\imp B_3,
\]
\[
(B_1\vee B_2)\imp B_3\equiv(B_1\imp B_3)\wedge(B_2\imp B_3),\qquad
B_1\imp\forall x\,B_2\equiv\forall x\,(B_1\imp B_2)\ (x\notin\operatorname{fv}(B_1)),
\]
together with distributivity of $\wedge$ and $\vee$ and the others listed on \book{64}.

\begin{example}[Size and behaviour {\book{64--66}}]\label{ex:horn-size}
In the common presentation, $((p\vee r)\wedge(q\vee t))\imp s$ must become four clauses
\[
(p\wedge q)\imp s,\quad (r\wedge q)\imp s,\quad (p\wedge t)\imp s,\quad (r\wedge t)\imp s,
\]
which duplicates $p, q, r, t$, exponentially in general. It also changes behaviour: with the original
clause $p$ is derived at most once per use, with the four clauses it is derived again for each clause
that mentions it. Introducing new predicates $pr, qt$ for $p\vee r$ and $q\vee t$ instead,
\[
p\imp pr,\quad r\imp pr,\quad q\imp qt,\quad t\imp qt,\quad (pr\wedge qt)\imp s,
\]
grows the program only linearly. This is not a logical equivalence (there are new constants), but it is
\emph{conservative}: for every goal $G$ over the original signature $\Sig_1$,
$\seq{\Sig_1}{\Prog,D_1}{G}$ is provable iff $\seq{\Sig_2}{\Prog,D_2}{G}$ is.
\end{example}

\section{The role of types}\label{sec:horn-types}

\paragraph{Types classify expressions.} In the first-order setting the only distinction types
\emph{force} is between propositions (\lp{o}) and everything else \book{66}. Equality is syntactic: ``two
closed expressions are equal only if they are exactly the same'', so \lp{2 + 3}, \lp{3 + 2} and \lp{5}
are different elements of \lp{int}, and \lp{o} is a set of formulas, not of truth values
\book{66--67}.

\paragraph{Polymorphism, transparency, determinacy.} The declaration \lp{type nil list A.} is read as
$\forall A\,(\mathit{type}\ \mathit{nil}\ (\mathit{list}\ A))$: \lp{nil} has many types \book{67}.
\begin{definition}[Transparent, determinate {\book{67--68}}]
For an order-1 type $\tau_1\to\cdots\to\tau_n\to\tau_0$, a type variable occurring in $\tau_0$ is
\emph{transparent}; the type is \emph{determinate} if all its type variables are transparent.
\end{definition}
Knowing the target type of a determinate constructor fixes the types of its arguments: all elements of a
\lp{list A} have the same type. The constructor \lp{type cons A -> lst -> lst.} is not determinate and
builds heterogeneous lists \book{68}.

\paragraph{Checking and inference.} λProlog is strongly typed: types of constants are declared, types of
variables are inferred most generally, and checking is static \book{68--69}. The asymmetry of
\cref{sec:horn-terms} returns: \lp{append (1::nil) (2::nil) X, append ("abc"::nil) ("efg"::nil) Y.} is well
typed, with the two \lp{append} occurrences at different instances, but replacing \lp{Y} by \lp{X} makes it
ill typed, since a variable has one type \book{68}.

\paragraph{Types at run time.} With polymorphism and non-determinate constructors, types cannot always be
compiled away \book{69}. The book's Figure~2.8 builds heterogeneous lists with the non-determinate
\lp{cons} and dispatches on the run-time type of their elements \book{70}:
\begin{lprolog}
kind lst       type.
type null      lst.
type cons      A -> lst -> lst.
type separate  lst -> list int -> list real -> o.
separate (cons (X:int)  L) (X::K) M :- separate L K M.
separate (cons (X:real) L) K (X::M) :- separate L K M.
separate null nil nil.
\end{lprolog}
The annotation \lp{(X:int)} must be matched, during the search, against the type of the head element of
the incoming list: this is run-time type processing. Figure~2.9 achieves the same with
\emph{injections}, so that the dispatch is on \emph{terms} \book{70}:
\begin{lprolog}
kind numb      type.
type inj_int   int  -> numb.
type inj_real  real -> numb.
type separate  list numb -> list int -> list real -> o.
separate ((inj_int  X)::L) (X::K) M :- separate L K M.
separate ((inj_real X)::L) K (X::M) :- separate L K M.
separate nil nil nil.
\end{lprolog}
The book recommends determinate constructors and injections \book{71--72}: static typing then catches more
errors, and less type information has to be carried during proof search. Giving constants the type \lp{A}
would make nonsense well typed. The book poses the queries
\begin{lprolog}
separate (cons 1.0 (cons 2 (cons 3.0 null))) L K
separate ((inj_real 1.0)::(inj_int 2)::(inj_real 3.0)::nil) L K
\end{lprolog}
to the two figures and says that both bind ``K to (2 :: nil) and L to (1.0 :: 3.0 :: nil)'' \book{70}.
That swaps the answers: the second argument collects the integers, so \lp{L = (2 :: nil)} and
\lp{K = (1.0 :: 3.0 :: nil)}. The \micro{} port below confirms it.

\begin{kanren}[Types and \micro]
\micro{} has one universe of Scheme data. The book observes that in the first-order setting the only
distinction types \emph{force} is between \lp{o} and the rest, so all other sorts could be collapsed into
one \book{66}; \micro{} is that collapsed case. What the collapse gives up is static error detection,
typed unification (U7 above) and the run-time dispatch of Figure~2.8. Injections (Figure~2.9) are ordinary
miniKanren style, tagged lists such as \scm{(inj_int 2)}, and port clause by clause (file
\texttt{separate.scm}):
\begin{scheme}
(define-relation (separate° ns is rs)
  (cond° ((fresh° (x l k) (=° ns `((inj_int ,x) . ,l)) (=° is (cons x k)) (separate° l k rs)))
         ((fresh° (x l m) (=° ns `((inj_real ,x) . ,l)) (=° rs (cons x m)) (separate° l is m)))
         ((=° ns '()) (=° is '()) (=° rs '()))))

(μkanren-run (q -1 #t)
  (fresh° (l k) (=° q (list l k))
    (separate° '((inj_real 1.0) (inj_int 2) (inj_real 3.0)) l k)))
\end{scheme}
\begin{answers}
separate (...) L K => (((2) (1.0 3.0)))
\end{answers}
The answer is the pair $(L, K)$: the integers first, as the types demand. Where a type layer could live (a static
checker written as a relation over declared signatures, since type unification is just \scm{=°} on type
terms) is left open: the proposal stays untyped (\cref{sec:roadmap-lim-untyped}); simple types are a later
work item (\cref{sec:roadmap-later-types}).
\end{kanren}

\section{From fohc to \texorpdfstring{\texttt{(aux kanren micro)}}{(aux kanren micro)}}\label{sec:horn-micro}

We can now state precisely how \fohc{} maps onto \micro{} (described in \cref{ch:micro}):

\begin{center}
\small
\begin{tabular}{@{}>{\raggedright\arraybackslash}p{0.27\linewidth}>{\raggedright\arraybackslash}p{0.33\linewidth}>{\raggedright\arraybackslash}p{0.33\linewidth}@{}}
\toprule
\fohc{} (book) & \micro & remark\\
\midrule
$\top$, \texttt{true} ($\top R$) & \texttt{✓°} & \\
$G_1\wedge G_2$, \texttt{,} ($\wedge R$) & \texttt{(and° g1 g2)} & bind interleaves answers\\
$G_1\vee G_2$, \texttt{;} ($\vee R$) & \texttt{(or° g1 g2)}, \texttt{cond°} & interleaving, not left first\\
$\exists x\,G$, \texttt{sigma} ($\exists R$, logic variable) & \texttt{(fresh° (x) g)} & exactly the rule of \book{61}\\
atomic goal, \rulename{decide} & call \texttt{(p° t ...)} & clause choice compiled into the relation\\
clauses of $p$, $\forall\bar x\,(G\imp p\,\bar t)$ & \texttt{define-relation} with one \texttt{cond°} line per clause & precompiled and indexed by name, \book{61--62}\\
$\forall L$ with logic variable & \texttt{fresh°} opening each line & renaming apart is automatic\\
\rulename{initial} with unification & \texttt{=°} against the clause head & the threaded substitution percolates $\theta$\\
occurs check & built in & as the book requires, \book{32}\\
answer substitution & reified answer, \texttt{α}, \texttt{β}, \dots & up to renaming\\
typed unification & none & see U7\\
\texttt{<}, \texttt{>=}, \texttt{is} & \texttt{project°} + Scheme & non-logical, moded\\
\texttt{!} (pruning) & no hard cut; soft cut \texttt{if°}, \texttt{take°} & control, not logic\\
\texttt{not} & \texttt{(if° g ✗° ✓°)} & same unsoundness\\
${\imp}R$ (\rulename{augment}), $\forall R$ (\rulename{generic}) & absent & not needed in \fohc{}; \cref{ch:hh}\\
\bottomrule
\end{tabular}
\end{center}

So pure miniKanren \emph{is} \fohc, with two caveats. First, the shape of a relation (one disjunction of
existentially quantified equations per predicate) is the precompiled form of \book{61--62}, and the
built-in \scm{=°} is definable in \fohc{} by the single clause $\forall x\,(\cst{eq}\ x\ x)$.
Second, a Scheme-hosted miniKanren is strictly \emph{larger} than \fohc: goals are first-class Scheme
values, and \scm{project°} leaves the logic altogether. The \fohc{} fragment is a discipline, not something
the host enforces.

\begin{kanren}[The FSM, side by side with its λProlog original]
File \texttt{fsm.scm} ports Figures~2.4 and~2.5 clause by clause (\scm{append°} as above):
\begin{scheme}
;; Figure 2.4
(define-relation (path° s letters t)
  (cond° ((=° letters '()) (=° s t))
         ((fresh° (arc m rest)
            (trans° s arc m) (append° arc rest letters) (path° m rest t)))))
(define-relation (accept° w)
  (fresh° (s f) (start° s) (path° s w f) (final° f)))

;; Figure 2.5
(define-relation (start° s) (=° s 'q1))
(define-relation (final° s) (cond° ((=° s 'q2)) ((=° s 'q3))))
(define-relation (trans° s arc t)
  (cond° ((=° (list s arc t) '(q1 (a) q1)))
         ((=° (list s arc t) '(q1 (b) q1)))
         ((=° (list s arc t) '(q1 (a b) q2)))
         ((=° (list s arc t) '(q1 (b a) q3)))))

(define-relation (lists° l)
  (cond° ((=° l '())) ((fresh° (h t) (=° l (cons h t)) (lists° t)))))

(μkanren-run (q -1 #t) (accept° '(b b a b)))
(μkanren-run (q -1 #t) (fresh° (x y) (=° q (list x y)) (accept° `(b a ,x ,y))))
(μkanren-run (q 4 #t) (lists° q))
(μkanren-run (q 5 #t) (lists° q) (accept° q))
(μkanren-run (q 5 #t) (accept° q))
\end{scheme}
\begin{answers}
accept (b b a b)         => (α)
accept (b a X Y)         => ((a b) (b a))
lists L          (4)     => (() (α) (α β) (α β γ))
lists L, accept L (5)    => ((a b) (b a) (a a b) (a b a) (b a b))
accept L          (5)    => ((a b) (b a) (a a b) (a b a) (b a b))
\end{answers}
The first four lines agree with the book (\lp{lists L} gives \lp{nil}, \lp{T1 :: nil}, \dots). The last
line is the interesting one: \scm{(accept° q)} \emph{without} the generator produces answers, where
depth-first search loops on the $a$-arc. The order happens to coincide with the book's here; in general
one should not expect \micro{} to reproduce λProlog's answer order.
\end{kanren}

\subsection{Depth-first search versus interleaving}\label{sec:horn-dfs}

λProlog commits to the depth-first, left-to-right, clause-order strategy of \cref{sec:horn-pragmatics}.
Teyjus compiles it to an abstract machine \book{277}. \micro{} keeps the same goals but changes the
strategy, as μKanren does~\cite{HemannFriedman2013}: \scm{or°} interleaves its two answer streams, switching
on each suspension, and a relation call suspends. A clause that recurses forever therefore gets only a
share of the search, and the remaining clauses still run. That is why \scm{(accept° q)} terminates
answer by answer.

What does this buy, logically? \Cref{thm:horn-oproof} is a statement about the proof system: O-proofs exist
whenever proofs exist. A search strategy is complete if it finds a proof whenever one exists, and, for answers, if
every answer substitution of an O-proof is eventually produced (up to instantiation). Depth-first search
is not \book{60}. For miniKanren's interleaving search, Rozplokhas, Vyatkin and
Boulytchev~\cite{RozplokhasVyatkinBoulytchev2020} proved soundness and completeness with respect to a
denotational semantics: every answer is eventually produced. \micro's particular interleaving has not
been checked against the hypotheses of that proof, so read it as a strong expectation for \micro, not a
theorem about it. Interleaving is fair for \emph{disjunction} only. Termination is not
promised: a query with no answers, or the rest of a query after its last answer, can still diverge,
which is unavoidable since \fohc{} provability is undecidable \book{51}. Conjunction remains
order-sensitive. For example
\begin{scheme}
(μkanren-run (q 1 #t) (lists° q) (=° q 'not-a-list))
\end{scheme}
never returns: \scm{lists°} offers infinitely many lists, and each is rejected by the second conjunct. We
ran it with a 10-second timeout, and it produced no output. With the conjuncts swapped it fails at once.

\begin{pitfall}[Answer order is not part of the logic]
The book stresses that programs are lists, so the order and multiplicity of clauses are observable under
depth-first search \book{60}, and logically equivalent transformations can change the number of
derivations (\cref{ex:horn-size}). Under interleaving the same holds, with a different order. A port of a
book example to \micro{} should compare answer \emph{sets}, unless the order has been checked.
\end{pitfall}

\begin{kanren}[Negation as failure and moded built-ins in \micro]
\micro{} has a soft cut, \scm{(if° g? gt gf)}: if \scm{g?} has answers, all of them continue with \scm{gt},
otherwise \scm{gf} runs. With it, \lp{not} is one line, and the book's commutativity failure \book{137--138}
reproduces exactly (file \texttt{naf.scm}). For contrast, the last two lines use the disequality constraint
\scm{≠°}, which is not part of \fohc{} (\cref{ch:micro}) but is commutative:
\begin{scheme}
(define (not° g) (if° g ✗° ✓°))

(μkanren-run (q -1 #t) (=° q 2) (not° (=° 1 q)))
(μkanren-run (q -1 #t) (not° (=° 1 q)) (=° q 2))
(μkanren-run (q -1 #t) (=° q 2) (≠° 1 q))
(μkanren-run (q -1 #t) (≠° 1 q) (=° q 2))
\end{scheme}
\begin{answers}
X = 2, not (1 = X)   => (2)
not (1 = X), X = 2   => ()
X = 2, X =/= 1       => (2)
X =/= 1, X = 2       => (2)
\end{answers}
Moded built-ins behave as in λProlog. The port of \lp{insert} guards with \scm{project°}, which hands the
current values to Scheme (file \texttt{insert.scm}):
\begin{scheme}
(define-relation (insert° x t t*)
  (cond° ((=° t 'empty) (=° t* `(node ,x empty empty)))
         ((fresh° (a l r nl)
            (=° t `(node ,a ,l ,r)) (=° t* `(node ,a ,nl ,r))
            (project° ((x° x) (a° a)) (if (< x° a°) ✓° ✗°))
            (insert° x l nl)))
         ((fresh° (a l r nr)
            (=° t `(node ,a ,l ,r)) (=° t* `(node ,a ,l ,nr))
            (project° ((x° x) (a° a)) (if (>= x° a°) ✓° ✗°))
            (insert° x r nr)))))

(μkanren-run (q -1 #t) (insert° 4 '(node 3 (node 2 empty empty) empty) q))
(μkanren-run (q -1 #t) (fresh° (x) (insert° x '(node 3 empty empty) q)))   ; mode broken
\end{scheme}
\begin{answers}
insert 4 (node 3 (node 2 empty empty) empty) T => ((node 3 (node 2 empty empty) (node 4 empty empty)))
insert X (node 3 empty empty) T => 
Error: (<) bad argument type - not a number: _1
\end{answers}
The second query breaks the mode: Scheme's \scm{<} receives the unbound logic variable \scm{_1} and the run
aborts with the error shown (printed on standard error). Likewise intensional
arithmetic (file \texttt{arith.scm}):
\begin{scheme}
(μkanren-run (q -1 #t) (=° '(+ 2 3) 5))
(μkanren-run (q -1 #t) (=° q '(+ 2 3)))
(μkanren-run (q -1 #t)
  (fresh° (e) (=° e '(+ 2 3)) (project° ((e* e)) (=° q (apply + (cdr e*))))))
\end{scheme}
\begin{answers}
2 + 3 = 5        => ()
X = 2 + 3        => ((+ 2 3))
X is 2 + 3       => (5)
\end{answers}
\end{kanren}

\begin{kanren}[What \textup{fohc} does not need, and later chapters will]
By \cref{prop:horn-flat}, nothing in this chapter needs the program or the signature to change during
the search, so \micro{} can keep relations as closed Scheme procedures and needs no signature object.
\Cref{ch:hh} adds \rulename{augment} (clauses that are visible only along one branch of the search) and
\rulename{generic} (fresh constants, together with a scoping check in unification that the occurs-check
discussion of \book{32} anticipates). \Cref{ch:lambda,ch:hounif} add $\lambda$-terms as data. The
corresponding requirements on \micro{} are collected in \cref{ch:design}.
\end{kanren}

\section{Summary}\label{sec:horn-summary}
\begin{itemize}
  \item Logic programming is \emph{computation as proof search}: the state is a sequent
    $\seq{\Sig}{\Prog}{G}$, computing means looking for a \emph{cut-free} proof, and cut is used to reason
    about programs, not to run them \book{3,63}.
  \item Data are typed first-order terms. Kinds are arities. Each occurrence of a polymorphic constant
    takes its own instance of the type, but a variable has one type. Principal types exist. First-order
    encodings of syntax do not capture binding \book{10--25}.
  \item First-order unification is a terminating transformation system on multisets of equations. It computes
    mgu's, unique up to renaming. The occurs check is mandatory in the book. Types matter only when two
    heads have the same name, and they can be erased for parametric signatures \book{26--32}.
  \item Connectives in goals have a fixed search semantics, the right rules (\rulename{and}, \rulename{or},
    \rulename{instan}, \rulename{augment}, \rulename{generic}, \rulename{true}); atoms are solved by
    \emph{backchaining}, the focused left rules \rulename{decide}, \rulename{initial}, ${\imp}L$, $\wedge L$,
    $\forall L$ \book{39--45}.
  \item In \fohc{} ($G ::= \top\mid A\mid G\wedge G\mid G\vee G\mid\exists x\,G$; $D ::= A\mid G\imp D\mid D\wedge
    D\mid\forall x\,D$) O-proofs are complete for classical and intuitionistic logic, and the program and
    signature never change during the search \book{43--46}.
  \item λProlog fixes depth-first, left-to-right, clause-order search with logic variables and
    unification. It is predictable but incomplete. \lp{!}, \lp{not}, \lp{is} and comparisons are
    non-logical \book{60--62,137,286}.
  \item Logically equivalent presentations of clauses prove the same goals but may differ in size and in the
    number of derivations; introducing predicates is conservative \book{63--66}.
  \item Types classify expressions, not values; equality is syntactic. Determinate constructors and
    injections are preferred to run-time type dispatch \book{66--72}.
  \item Pure miniKanren is \fohc{}, without types and with an interleaving search. The interleaving search
    finds answers that depth-first search misses (\scm{(accept° q)}), but conjunction order still decides
    termination. \micro{} reproduces the book's examples and its warnings (ill-typed unifier,
    non-commutative \lp{not}, moded built-ins).
\end{itemize}
